%!TEX root = ejemplo.tex
%=========================================================
\section{Catálogo de mensajes}	
\label{sec:mensajes}

En esta sección se describen todos los mensajes que aparecen en el sistema. Para cada mensaje se especifica:

\begin{description}\itemsep0em
	\item[Id:] Identificador del mensaje de la forma ``MSG XX'' y descripción corta del mismo.
	\item[Tipo:] Tipo del mensaje el cual puede ser: 
	\begin{description}
		\item[Normal:] Mensaje que informa al usuario una instrucción o el estado interno que guarda el sistema, suele tener un color {\color{msgNormalColor}Azul}.
		\item[Éxito:] Mensaje que informa al usuario sobre una acción realizada, sirve para confirmar el correcto funcionamiento del sistema. Se presentan con un color {\color{msgInfoColor}Verde}.
		\item[Atención:] Mensaje que tiene como finalidad llamar la atención del usuario a una situación que requiere su intervención, por ejemplo cuando una actividad ha generado un efecto colateral o se realizará una acción destructiva y no reversible. Se presentan con un color {\color{msgWarningColor}Naranja}.
		\item[Error:] Mensaje que informa al usuario un fallo en en una operación o un impedimento para realizarla, por ejemplo: cuando no se puede efectuar la acción solicitada, cuando un dato falta o tiene un formato no aceptado por el sistema. Se presentan con un color {\color{msgErrorColor}Rojo}.
	\end{description}
	\item[Propósito:] Explicación del propósito del mensaje.
	\item[Redacción:] Redacción del mensaje.
	\item[Parámetros:] En caso de que el mensaje pueda variar se especifican los casos y la forma en que debe adaptarse la redacción
	\item[Ejemplos:] Ejemplos de como debe renderizarse el mensaje.
\end{description}


\subsection{Lista de mensajes}

%--------------------------------------
% MENSAJES DE CONTROL DE ACCESO
%--------------------------------------

\begin{cdtMessage}[msgInfoColor]{MSG01}{Bienvenida al usuario}
	\item[Propósito:] Indicar al usuario que ha ingresado 
	satisfactoriamente al sistema SICORRE.
	\item[Redacción:] Bienvenido/a $<$nombre$>$.
	\item[Parámetros:] \hspace{1cm}
	\begin{itemize}
		\item $<$nombre$>$ Nombre completo del usuario autenticado.
	\end{itemize}
	\item[Ejemplos:] Bienvenido/a Joji Pérez Ramírez.
\end{cdtMessage}

\begin{cdtMessage}[msgErrorColor]{MSG02}{Credenciales incorrectas}
	\item[Propósito:] Indicar al usuario que el correo electrónico 
	o la contraseña ingresados no son válidos.
	\item[Redacción:] El correo electrónico o la contraseña 
	son incorrectos. Verifique sus datos e intente de nuevo.
	\item[Parámetros:] \hspace{1cm}
	\begin{itemize}
		\item Ninguno. Por seguridad no se especifica cuál de 
		los dos campos es incorrecto.
	\end{itemize}
	\item[Ejemplos:] El correo electrónico o la contraseña 
	son incorrectos. Verifique sus datos e intente de nuevo.
\end{cdtMessage}

\begin{cdtMessage}[msgErrorColor]{MSG03}{Acceso revocado}
	\item[Propósito:] Indicar al usuario que su cuenta ha sido 
	revocada y no puede acceder al sistema.
	\item[Redacción:] El acceso de $<$nombre$>$ ha sido revocado. 
	Contacte al administrador para más información.
	\item[Parámetros:] \hspace{1cm}
	\begin{itemize}
		\item $<$nombre$>$ Nombre completo del usuario 
		que intenta acceder.
	\end{itemize}
	\item[Ejemplos:] El acceso de Joji Pérez Ramírez ha sido 
	revocado. Contacte al administrador para más información.
\end{cdtMessage}

%--------------------------------------
% MENSAJES DE GESTIÓN DE USUARIOS
%--------------------------------------

\begin{cdtMessage}[msgWarningColor]{MSG04}{Confirmar revocación de acceso}
	\item[Propósito:] Solicitar confirmación al administrador 
	antes de revocar el acceso de un usuario al sistema.
	\item[Redacción:] ¿Está seguro que desea revocar el acceso 
	de $<$nombre$>$? El usuario no podrá ingresar al sistema, 
	pero su historial de actividad se conservará.
	\item[Parámetros:] \hspace{1cm}
	\begin{itemize}
		\item $<$nombre$>$ Nombre completo del usuario 
		a quien se revocará el acceso.
	\end{itemize}
	\item[Ejemplos:] ¿Está seguro que desea revocar el acceso 
	de Joji Pérez Ramírez? El usuario no podrá ingresar al 
	sistema, pero su historial de actividad se conservará.
\end{cdtMessage}

\begin{cdtMessage}[msgWarningColor]{MSG05}{Confirmar eliminación de usuario}
	\item[Propósito:] Solicitar confirmación al administrador 
	antes de eliminar permanentemente el registro de un usuario.
	\item[Redacción:] ¿Está seguro que desea eliminar el registro 
	de $<$nombre$>$? Esta acción no puede deshacerse.
	\item[Parámetros:] \hspace{1cm}
	\begin{itemize}
		\item $<$nombre$>$ Nombre completo del usuario 
		a eliminar del sistema.
	\end{itemize}
	\item[Ejemplos:] ¿Está seguro que desea eliminar el registro 
	de Joji Pérez Ramírez? Esta acción no puede deshacerse.
\end{cdtMessage}

\begin{cdtMessage}[msgInfoColor]{MSG06}{Registro exitoso de usuario}
	\item[Propósito:] Confirmar al administrador que el nuevo 
	usuario ha sido registrado correctamente en el sistema.
	\item[Redacción:] El usuario $<$nombre$>$ ha sido registrado 
	exitosamente en el sistema.
	\item[Parámetros:] \hspace{1cm}
	\begin{itemize}
		\item $<$nombre$>$ Nombre completo del usuario 
		recién registrado.
	\end{itemize}
	\item[Ejemplos:] El usuario Joji Pérez Ramírez ha sido 
	registrado exitosamente en el sistema.
\end{cdtMessage}

\begin{cdtMessage}[msgErrorColor]{MSG07}{CURP ya registrada}
	\item[Propósito:] Indicar al administrador que la CURP 
	ingresada ya está asociada a otro usuario en el sistema.
	\item[Redacción:] La CURP $<$curp$>$ ya se encuentra 
	registrada en el sistema.
	\item[Parámetros:] \hspace{1cm}
	\begin{itemize}
		\item $<$curp$>$ CURP ingresada en el formulario.
	\end{itemize}
	\item[Ejemplos:] La CURP PERJ900101HDFRZN09 ya se encuentra 
	registrada en el sistema.
\end{cdtMessage}

\begin{cdtMessage}[msgErrorColor]{MSG08}{Correo ya registrado}
	\item[Propósito:] Indicar al administrador que el correo 
	electrónico ingresado ya está asociado a otro usuario 
	en el sistema.
	\item[Redacción:] El correo $<$correo$>$ ya se encuentra 
	registrado en el sistema.
	\item[Parámetros:] \hspace{1cm}
	\begin{itemize}
		\item $<$correo$>$ Correo electrónico ingresado 
		en el formulario.
	\end{itemize}
	\item[Ejemplos:] El correo joji.perez@fundacion.org ya 
	se encuentra registrado en el sistema.
\end{cdtMessage}

\begin{cdtMessage}[msgErrorColor]{MSG09}{Código postal no encontrado}
	\item[Propósito:] Indicar al usuario que el código postal 
	ingresado no existe en el catálogo SEPOMEX del sistema.
	\item[Redacción:] El código postal $<$cp$>$ no fue encontrado. 
	Verifique e intente de nuevo.
	\item[Parámetros:] \hspace{1cm}
	\begin{itemize}
		\item $<$cp$>$ Código postal ingresado en el formulario.
	\end{itemize}
	\item[Ejemplos:] El código postal 00000 no fue encontrado. 
	Verifique e intente de nuevo.
\end{cdtMessage}

\begin{cdtMessage}[msgErrorColor]{MSG10}{Campos obligatorios incompletos}
	\item[Propósito:] Indicar al usuario que existen campos 
	obligatorios sin completar antes de guardar el formulario.
	\item[Redacción:] El campo $<$campo$>$ es obligatorio 
	y no puede estar vacío.
	\item[Parámetros:] \hspace{1cm}
	\begin{itemize}
		\item $<$campo$>$ Nombre del campo obligatorio 
		que se encuentra vacío.
	\end{itemize}
	\item[Ejemplos:] El campo Nombre(s) es obligatorio 
	y no puede estar vacío.
\end{cdtMessage}

\begin{cdtMessage}[msgInfoColor]{MSG11}{Modificación exitosa de usuario}
	\item[Propósito:] Confirmar al administrador que los datos 
	del usuario han sido actualizados correctamente en el sistema.
	\item[Redacción:] Los datos de $<$nombre$>$ han sido 
	actualizados exitosamente.
	\item[Parámetros:] \hspace{1cm}
	\begin{itemize}
		\item $<$nombre$>$ Nombre completo del usuario modificado.
	\end{itemize}
	\item[Ejemplos:] Los datos de Joji Pérez Ramírez han sido 
	actualizados exitosamente.
\end{cdtMessage}

%--------------------------------------
% MENSAJES DE GESTIÓN DE EQUIPOS
%--------------------------------------

\begin{cdtMessage}[msgWarningColor]{MSG12}{Sin equipos registrados}
	\item[Propósito:] Informar al administrador que no existe 
	ningún equipo multidisciplinario registrado en el sistema.
	\item[Redacción:] No hay equipos multidisciplinarios 
	registrados en el sistema.
	\item[Parámetros:] \hspace{1cm}
	\begin{itemize}
		\item Ninguno.
	\end{itemize}
	\item[Ejemplos:] No hay equipos multidisciplinarios 
	registrados en el sistema.
\end{cdtMessage}

\begin{cdtMessage}[msgWarningColor]{MSG13}{Confirmar eliminación de equipo}
	\item[Propósito:] Solicitar confirmación al administrador 
	antes de eliminar permanentemente el registro de un equipo 
	multidisciplinario.
	\item[Redacción:] ¿Está seguro que desea eliminar el equipo 
	$<$nombre$>$? Esta acción no puede deshacerse.
	\item[Parámetros:] \hspace{1cm}
	\begin{itemize}
		\item $<$nombre$>$ Nombre del equipo a eliminar.
	\end{itemize}
	\item[Ejemplos:] ¿Está seguro que desea eliminar el equipo 
	Equipo Alpha? Esta acción no puede deshacerse.
\end{cdtMessage}

\begin{cdtMessage}[msgInfoColor]{MSG14}{Equipo registrado exitosamente}
	\item[Propósito:] Confirmar al administrador que el nuevo 
	equipo multidisciplinario ha sido registrado correctamente 
	en el sistema.
	\item[Redacción:] El equipo $<$nombre$>$ ha sido registrado 
	exitosamente en el sistema.
	\item[Parámetros:] \hspace{1cm}
	\begin{itemize}
		\item $<$nombre$>$ Nombre del equipo recién registrado.
	\end{itemize}
	\item[Ejemplos:] El equipo Equipo Alpha ha sido registrado 
	exitosamente en el sistema.
\end{cdtMessage}

\begin{cdtMessage}[msgErrorColor]{MSG15}{Nombre de equipo ya registrado}
	\item[Propósito:] Indicar al administrador que el nombre 
	ingresado para el equipo ya existe en el sistema.
	\item[Redacción:] El nombre de equipo $<$nombre$>$ ya se 
	encuentra registrado en el sistema.
	\item[Parámetros:] \hspace{1cm}
	\begin{itemize}
		\item $<$nombre$>$ Nombre de equipo ingresado 
		en el formulario.
	\end{itemize}
	\item[Ejemplos:] El nombre de equipo Equipo Alpha ya se 
	encuentra registrado en el sistema.
\end{cdtMessage}

\begin{cdtMessage}[msgWarningColor]{MSG16}{No hay integrantes disponibles}
	\item[Propósito:] Informar al administrador que todo el 
	personal activo ya pertenece a un equipo multidisciplinario 
	y no hay integrantes disponibles para asignar.
	\item[Redacción:] No hay personal disponible para agregar 
	al equipo. Todos los integrantes activos ya pertenecen 
	a un equipo multidisciplinario.
	\item[Parámetros:] \hspace{1cm}
	\begin{itemize}
		\item Ninguno.
	\end{itemize}
	\item[Ejemplos:] No hay personal disponible para agregar 
	al equipo. Todos los integrantes activos ya pertenecen 
	a un equipo multidisciplinario.
\end{cdtMessage}

\begin{cdtMessage}[msgInfoColor]{MSG17}{Integrante agregado exitosamente}
	\item[Propósito:] Confirmar al administrador que el integrante 
	seleccionado ha sido asignado correctamente al equipo 
	multidisciplinario.
	\item[Redacción:] $<$nombre$>$ ha sido agregado exitosamente 
	al equipo $<$equipo$>$.
	\item[Parámetros:] \hspace{1cm}
	\begin{itemize}
		\item $<$nombre$>$ Nombre completo del integrante agregado.
		\item $<$equipo$>$ Nombre del equipo al que fue asignado.
	\end{itemize}
	\item[Ejemplos:] Joji Pérez Ramírez ha sido agregado 
	exitosamente al equipo Equipo Alpha.
\end{cdtMessage}

%--------------------------------------
% MENSAJES DE CONSULTA GENERAL
%--------------------------------------

\begin{cdtMessage}[msgWarningColor]{MSG18}{Sin usuarios registrados}
	\item[Propósito:] Informar al administrador que no existe 
	ningún usuario registrado en el sistema.
	\item[Redacción:] No hay usuarios registrados en el sistema.
	\item[Parámetros:] \hspace{1cm}
	\begin{itemize}
		\item Ninguno.
	\end{itemize}
	\item[Ejemplos:] No hay usuarios registrados en el sistema.
\end{cdtMessage}
\begin{cdtMessage}[msgInfoColor]{MSG19}{Equipo modificado exitosamente}
	\begin{description}\itemsep0em
		\item[Descripción:] El nombre del equipo fue actualizado correctamente en el sistema.
	\end{description}
\end{cdtMessage}

\begin{cdtMessage}[msgInfoColor]{MSG20}{Integrante eliminado exitosamente}
	\begin{description}\itemsep0em
		\item[Descripción:] El integrante fue removido del equipo correctamente.
	\end{description}
\end{cdtMessage}

\begin{cdtMessage}[msgWarningColor]{MSG21}{No hay NNAs registrados}
	\begin{description}\itemsep0em
		\item[Descripción:] No existen NNAs registrados en el sistema.
	\end{description}
\end{cdtMessage}

\begin{cdtMessage}[msgWarningColor]{MSG22}{No hay tutores registrados}
	\begin{description}\itemsep0em
		\item[Descripción:] No existen tutores registrados en el sistema.
	\end{description}
\end{cdtMessage}

\begin{cdtMessage}[msgInfoColor]{MSG23}{NNA registrado exitosamente}
	\begin{description}\itemsep0em
		\item[Descripción:] El registro del NNA fue almacenado correctamente en el sistema.
	\end{description}
\end{cdtMessage}

\begin{cdtMessage}[msgErrorColor]{MSG24}{Error al registrar NNA}
	\begin{description}\itemsep0em
		\item[Descripción:] Ocurrió un problema al intentar guardar el registro del NNA.
	\end{description}
\end{cdtMessage}

\begin{cdtMessage}[msgInfoColor]{MSG25}{Tutor registrado exitosamente}
	\begin{description}\itemsep0em
		\item[Descripción:] El registro del tutor fue almacenado correctamente en el sistema.
	\end{description}
\end{cdtMessage}

\begin{cdtMessage}[msgErrorColor]{MSG26}{Error al registrar tutor}
	\begin{description}\itemsep0em
		\item[Descripción:] Ocurrió un problema al intentar guardar el registro del tutor.
	\end{description}
\end{cdtMessage}

\begin{cdtMessage}[msgWarningColor]{MSG27}{Sin registros disponibles}
	\begin{description}\itemsep0em
		\item[Descripción:] No hay registros disponibles para mostrar en el catálogo.
	\end{description}
\end{cdtMessage}